%%
%% Blank Wiley (WileyNJDv5) article skeleton — full-feature version
%% Stripped of lorem-ipsum sample body text — structure only.
%%
%% Compiler note: as of Feb 2024, Wiley requires XeLaTeX (not pdfLaTeX)
%% for font compatibility, with TeX Live 2022 (TeX Live 2023 will not work).
%%
%% Documentclass options select citation style and column layout:
%%   HARVARD / AMA / VANCOUVER / NUMBERED -> citation/reference style
%%   LATO2COL (two-column) or Times1COL (one-column, Times font) etc.
%% Swap these per the target journal's author guidelines.
%%\documentclass[AMA,Times1COL]{WileyNJDv5}
\documentclass[HARVARD,LATO2COL]{WileyNJDv5}

%% ---------------------------------------------------------------
%% Journal/article metadata.
%% ---------------------------------------------------------------
\articletype{}%
\received{}
\revised{}
\accepted{}
\journal{}
\volume{}
\copyyear{}
\startpage{}

\raggedbottom

\begin{document}

\title{}

%% One \author[n]{} per author; the [n] list links to matching
%% \address[n]{} entries below for shared/multiple affiliations.
\author[1]{}
\author[2,3]{}
\author[3]{}

%% Running header short author list, e.g. "Smith et al."
\authormark{}
%% Running header short title
\titlemark{}

%% \address[n]{} — repeat per distinct affiliation, matching the
%% \author[n] numbering above.
\address[1]{\orgdiv{}, \orgname{}, \orgaddress{\state{}, \country{}}}
\address[2]{\orgdiv{}, \orgname{}, \orgaddress{\state{}, \country{}}}
\address[3]{\orgdiv{}, \orgname{}, \orgaddress{\state{}, \country{}}}

%% Corresponding author note + contact email.
\corres{\email{}}

%% Optional: address at time of publication if different from \address.
%%\presentaddress{}

%% Optional funding/classification metadata:
%%\fundingInfo{}
%%\JELinfo{}

%% Optional short-title argument in [ ] for running heads/TOC.
\abstract[]{}

\keywords{}

%% Optional: how this article should itself be cited (used on cover
%% page / running header of some journals).
%%\jnlcitation{\cname{%
%%\author{}}.
%%\ctitle{} \cjournal{} \cvol{}}.

\maketitle

%% Optional abbreviations footnote on the first page (unnumbered):
%%\renewcommand\thefootnote{}
%%\footnotetext{\textbf{Abbreviations:} }
%%\renewcommand\thefootnote{\fnsymbol{footnote}}
%%\setcounter{footnote}{1}

\section{Introduction}\label{sec:intro}

\section{}\label{sec:generic}

\subsection{}

%% Full-width figure (spans both columns in two-column layout):
\begin{figure*}[t]
\centerline{\includegraphics[width=\textwidth,height=9pc]{}}
\caption{\label{fig:1}}
\end{figure*}

%% Single-column figure:
%%\begin{figure*}
%%\centerline{\includegraphics[width=342pt,height=9pc]{}}
%%\caption{\label{fig:2}}
%%\end{figure*}

\subsection{}

%% Quote environment:
%%\begin{quote}
%%
%%\end{quote}

\section{}\label{sec:another}

\subsection{}
\subsection{}
\subsubsection{}
\paragraph{}
\subparagraph{}

%% ---------------------------------------------------------------
%% Boxed callout content, with or without a heading.
%% ---------------------------------------------------------------
%%\begin{boxwithhead}
%%{BOX 1\quad }
%%{\noindent }
%%\end{boxwithhead}

%%\begin{boxtext}
%%{\noindent }
%%\end{boxtext}

%% ---------------------------------------------------------------
%% Full-width table with spanned column headings, \tnote markers, and
%% \tablenotes for footnotes/source line.
%% ---------------------------------------------------------------
\begin{center}
\begin{table*}[!h]%
\caption{\label{tab:1}}
\begin{tabular*}{\textwidth}{@{\extracolsep\fill}lllll@{}}
\toprule
&\multicolumn{2}{@{}l}{\textbf{}} & \multicolumn{2}{@{}l}{\textbf{}} \\\cmidrule{2-3}\cmidrule{4-5}
\textbf{} & \textbf{} & \textbf{} & \textbf{} & \textbf{} \\
\midrule
 & & & & \\
\bottomrule
\end{tabular*}
\begin{tablenotes}
\item[$^{\rm a}$]
\item {\it Source}:
\end{tablenotes}
\end{table*}
\end{center}

%% Sideways table (for wide tables) — spans \textheight:
%%\begin{sidewaystable}
%%\caption{\label{tab:2}}
%%\begin{tabular*}{\textheight}{@{\extracolsep\fill}lllllllll@{\extracolsep\fill}}
%%\toprule
%%\midrule
%%\bottomrule
%%\end{tabular*}
%%\begin{tablenotes}
%%\item[*]
%%\end{tablenotes}
%%\end{sidewaystable}

%% ---------------------------------------------------------------
%% List styles: bulleted, described, and numbered (with nested
%% alphabetical/roman sub-levels), plus an unnumbered variant.
%% ---------------------------------------------------------------
%%\begin{itemize}
%%\item
%%\end{itemize}

%%\begin{description}
%%\item[]
%%\end{description}

%%\begin{enumerate}[1.]
%%\item
%%\begin{enumerate}[a.]
%%\item
%%\begin{enumerate}[i.]
%%\item
%%\end{enumerate}
%%\end{enumerate}
%%\end{enumerate}

%%\begin{enumerate}[]
%%\item
%%\end{enumerate}

\section{Examples for Enunciations}\label{sec:theorems}

\begin{theorem}\label{thm:1}

\end{theorem}

\begin{proposition}

\end{proposition}

\begin{definition}

\end{definition}

\begin{proof}

\end{proof}

%% To attribute a proof to a specific theorem:
%%\begin{proof}[Proof of Theorem~{\rm\ref{thm:1}}]
%%
%%\end{proof}

%% Sideways figure (for wide images) — spans \textheight:
%%\begin{sidewaysfigure}
%%\centerline{\includegraphics[width=542pt,height=9pc]{}}
%%\caption{\label{fig:3}}
%%\end{sidewaysfigure}

%% Algorithm environment (algorithmicx/algpseudocode):
\begin{algorithm}
\caption{\enskip}\label{algo:1}
\begin{algorithmic}
\State
\end{algorithmic}
\end{algorithm}

%% Numbered display equation with \nonumber on intermediate lines,
%% \label only on the line that should carry a number:
\begin{align}\label{eq:1}
 &  \nonumber\\
 &
\end{align}

\section{Conclusions}\label{sec:conclusions}

%% ---------------------------------------------------------------
%% Backmatter sections. Use \bmsection*{} for unnumbered headings
%% (Author contributions, Acknowledgments, Financial disclosure,
%% Conflict of interest, Supporting information) — do not use \section*.
%% ---------------------------------------------------------------
\bmsection*{Author contributions}

\bmsection*{Acknowledgments}

\bmsection*{Financial disclosure}

\bmsection*{Conflict of interest}

\bibliography{wileyNJD-AMA}

\bmsection*{Supporting information}

%% ---------------------------------------------------------------
%% Appendices. Use \appendix once, then \bmsection{}/\bmsubsection{}/
%% \bmsubsubsection{} (not \section/\subsection) for appendix headings.
%% ---------------------------------------------------------------
\appendix

\bmsection{}\label{app:1}

%% Program code listing (listings package) within an appendix:
%%\begin{lstlisting}[caption={},label=,basicstyle=\fontsize{8}{10}\selectfont\ttfamily]
%%
%%\end{lstlisting}

\bmsubsection{}\label{app:1a}

\bmsubsubsection{}\label{app:1a1}

%% Unnumbered figure inside an appendix (no \caption):
%%\begin{center}
%%\includegraphics[width=7pc,height=8pc]{}
%%\end{center}

\bmsection{}\label{app:2}

%% Standard (non-appendix) figure inside an appendix section:
%%\begin{figure}[b]
%%\centerline{\includegraphics[height=10pc,width=78mm]{}}
%%\caption{\label{fig:4}}
%%\end{figure}

\bmsubsection{}\label{app:2a}

%% Compact (non-tabular*) table, centered, for appendix use:
%%\begin{center}
%%\begin{tabular*}{250pt}{@{\extracolsep\fill}lcc@{\extracolsep\fill}}
%%\toprule
%%\midrule
%%\bottomrule
%%\end{tabular*}
%%\end{center}

%% Show all bibliography entries, cited or not; comment out to show
%% only cited entries:
%%\nocite{*}

%% ---------------------------------------------------------------
%% Author biography — check the target journal's guidelines for
%% whether this is required (usually only for review-type articles).
%% ---------------------------------------------------------------
\bmsection*{Author Biography}

\begin{biography}{\includegraphics[width=76pt,height=76pt]{}}{
{\textbf{}}}
\end{biography}

\end{document}
%% End of blank WileyNJD skeleton.
